% Development assembled from DESARROLLO_MATEMATICO, sections 38--40 and 50.
\subsubsection{Capacity index, vacancies, and induction of the tritic regime}
\label{vac-capacidad-trit}

The sequence of vacancies is linked to refinement through an
increment identity. To establish it, fix the phase origin
\(\theta=0\). Let
\[
 \rho=\log_{10}9=1-\delta,\qquad
 \mathcal C_t=\lfloor(t+1)\rho\rfloor\quad(t\geq0).
\]
The index \(\mathcal C_t\) expresses the decimal capacity of the refinement
under the normalization that compares a ternary block of \(6\) positions,
with \(729\) possibilities, and a decimal block of \(3\) positions, with
\(1000\) possibilities. Indeed,
\(\log_{1000}729=\log_{10}9=\rho\). The letter \(K\) is reserved for
the dodecaphase register to be constructed subsequently.

\begin{proposition}[Complementarity of capacity and vacancy]
\label{prop:capacidad-vacancia}
For \(t\geq1\),
\begin{equation}
 \mathcal C_t-\mathcal C_{t-1}=1-z_t(0).
 \label{eq:capacidad-vacancia}
\end{equation}
Consequently, if \(0\leq a<b\),
\[
 \mathcal C_b-\mathcal C_a+
 \sum_{t=a+1}^{b}z_t(0)=b-a.
\]
\end{proposition}
\begin{proof}
The irrationality of \(\delta\) gives
\[
 \mathcal C_t
 =t+1-\lceil(t+1)\delta\rceil
 =t-\lfloor(t+1)\delta\rfloor.
\]
Subtracting two consecutive terms yields
\eqref{eq:capacidad-vacancia}. The sum over the interval telescopes.
\end{proof}

Each transition increases the block counter by one. If
\(z_t=0\), capacity also increases; if \(z_t=1\), capacity is
retained while a transition is incorporated into the history. Here a vacancy
is a capacity-retention event with a determined position and
antecedent. The preceding balance preserves exactly the
length between acquired capacity and recorded vacancies. The observed phase
\(g_t^{\rm cap}=[\mathcal C_t]_9\) satisfies
\[
 g_t^{\rm cap}-g_{t-1}^{\rm cap}\equiv1-z_t\pmod9.
\]
This capacity phase and the chronological phase \([t]_9\) are two
different observations. A local capacity retention and a full turn
of the holonomy have distinct mechanisms, although both
allow the phase observation to be preserved while modifying the state.

The relation with TRIT can be followed at the vacancy positions themselves.
For \(j\geq1\), define
\[
 p_j=\left\lfloor\frac j\delta\right\rfloor,\qquad
 v_j=\mathcal C_{p_j},\qquad h_j=p_{j+1}-p_j.
\]
The inequality \(p_j\delta<j<(p_j+1)\delta\) identifies \(p_j\)
as the position of event \(j\). Since \(0<\delta<1\), it also gives
\(\lfloor(p_j+1)\delta\rfloor=j\), and therefore
\[
 v_j=p_j-j,\qquad
 v_{j+1}-v_j=h_j-1.
\]
If \(s_j=22-h_j\), then \(s_j\in\{0,1\}\) and
\begin{equation}
 [v_{j+1}]_3=[v_j-s_j]_3.
 \label{eq:vacancia-regimen}
\end{equation}
An interval of length \(22\) preserves the class; one of length
\(21\) shifts it by one unit in the negative direction. The canonical selector
is applied to the positive phase \(r_j=1+[v_j]_9\):
\[
 \tau_j=M_{\rm ph}(r_j).
\]
A tritic regime is thereby determined by the capacity index
retained at each event, in addition to the binary coding of the vacancy.

The positions of the short intervals form the induced mechanical coding
of slope \(22-1/\delta\). Their separations \(6\) and \(7\)
determine the lengths of the plateaus of the class \([v_j]_3\), once
the origin and endpoint convention have been fixed. The first classes
are \(2\) for \(6\) events, \(1\) for \(7\), and \(0\) for
\(7\); the initial run of \(6\) is the initial restriction of a plateau.
Their signatures run through \(0,-1,+1\). Equations
\eqref{eq:capacidad-vacancia} and \eqref{eq:vacancia-regimen} make explicit the
link among refinement, vacancies, and regime selection. The quadratic
TRIT reading then realizes those types through their respective
generators; the temporal event and its geometric representation
thus preserve an identifiable dependency.

Inverting the monotone index yields another useful observation:
\[
 \mathcal T(k)=\min\{t:\mathcal C_t=k\}
   =\left\lceil\frac k\rho\right\rceil-1\qquad(k\geq1).
\]
If \(\sigma=\rho^{-1}-1\), the number of blocks required to gain
\(m\) units of capacity is
\[
 \mathcal T(k+m)-\mathcal T(k)
 =m+\lceil(k+m)\sigma\rceil-\lceil k\sigma\rceil
 \in\{m+\lfloor m\sigma\rfloor,m+\lceil m\sigma\rceil\}.
\]
In particular, \(9\) units of capacity require \(9\) or \(10\)
blocks. The period of a finite phase remains exact, whereas the
return time measured by the other index admits two values. This
distinction is part of the refinement dynamics and is not an
a posteriori correction of its numerical values.
